\documentclass[aspectratio=169]{beamer}

\usetheme{Madrid}
\usecolortheme{default}

% Clean, modern look
\setbeamertemplate{navigation symbols}{}
\setbeamertemplate{itemize items}[circle]

% Custom footline with name, institution, and page number
\setbeamertemplate{footline}{
  \leavevmode%
  \hbox{%
  \begin{beamercolorbox}[wd=.333333\paperwidth,ht=2.25ex,dp=1ex,left]{author in head/foot}%
    \hspace*{2ex}Luca J. Bolz
  \end{beamercolorbox}%
  \begin{beamercolorbox}[wd=.333333\paperwidth,ht=2.25ex,dp=1ex,center]{title in head/foot}%
    MPI Multidisciplinary Sciences
  \end{beamercolorbox}%
  \begin{beamercolorbox}[wd=.333333\paperwidth,ht=2.25ex,dp=1ex,right]{date in head/foot}%
    \insertframenumber{} / \inserttotalframenumber\hspace*{2ex}
  \end{beamercolorbox}}%
  \vskip0pt%
}

% Packages
\usepackage{booktabs}
\usepackage{graphicx}
\usepackage{amsmath}
\usepackage{xcolor}

% Colors
\definecolor{mpigreen}{RGB}{0,100,60}
\setbeamercolor{structure}{fg=mpigreen}
\setbeamercolor{title}{fg=white,bg=mpigreen}
\setbeamercolor{frametitle}{fg=white,bg=mpigreen}
\setbeamercolor{framesubtitle}{fg=white}

% Title info
\title{Machine Learning for MINFLUX Distance Estimation}
\subtitle{Fast, Accurate Distance Estimation with XGBoost}
\author{Luca J. Bolz}
\institute{
    Max Planck Institute for Multidisciplinary Sciences\\
    Kiel University
}
\date{15th of January 2026}

% Logo (add your MPI logo file as mpi_logo.png or mpi_logo.pdf)
% Uncomment when logo is available:
% \logo{\includegraphics[height=0.8cm]{mpi_logo.png}}

\begin{document}

% ============================================
% SLIDE 1: Title
% ============================================
\begin{frame}
    \titlepage
\end{frame}

% ============================================
% SLIDE 2: What We Measure
% ============================================
\begin{frame}{What We Measure}

    \textbf{Goal:} Distance $d$ between two fluorophores

    \vspace{0.8cm}

    \centering
    \begin{Large}
        [Fluorophore A] \hspace{0.5cm} $\xleftrightarrow{\hspace{1cm} d \hspace{1cm}}$ \hspace{0.5cm} [Fluorophore B]
    \end{Large}

    \vspace{1cm}

    \raggedright
    \textbf{How:}
    \begin{itemize}
        \item Excitation beam with intensity minimum (donut)
        \item Scan 6 positions: $X^-$, $X^0$, $X^+$, $Y^-$, $Y^0$, $Y^+$
        \item Photon count at minimum $\rightarrow$ position information
        \item Less signal = closer to minimum = better localization
    \end{itemize}

\end{frame}

% ============================================
% SLIDE 3: The Problem
% ============================================
\begin{frame}{The MLE Bottleneck}

    \textbf{Current approach: Maximum Likelihood Estimation (MLE)}
    \begin{itemize}
        \item Iterative: guess position $\rightarrow$ compute likelihood $\rightarrow$ update $\rightarrow$ repeat
        \item Requires 10--100 iterations until convergence
        \item $\sim$1--10 ms per estimation (implementation-dependent)
    \end{itemize}

    \vspace{0.8cm}

    \textbf{Problem:} Too slow for real-time tracking of molecular dynamics

    \vspace{0.5cm}

    \textbf{Goal:} Enable real-time distance estimation with comparable accuracy

\end{frame}

% ============================================
% SLIDE 3: The Solution
% ============================================
\begin{frame}{The Solution: XGBoost}

    \textbf{What is XGBoost?}
    \begin{itemize}
        \item Ensemble of 500 decision trees
        \item Each tree: simple if/else rules on features
        \item \textbf{Boosting:} each tree corrects errors of previous trees
        \item Final prediction = sum of all tree outputs
    \end{itemize}

    \vspace{0.5cm}

    \textbf{Why XGBoost?}
    \begin{itemize}
        \item Fast inference: just 500 $\times$ (compare $\rightarrow$ go left/right)
        \item No iterative optimization needed
        \item Learns non-linear relationships from data
    \end{itemize}

    \vspace{0.5cm}

    \textbf{Training:}
    \begin{itemize}
        \item 75k samples (balanced: 25k per distance class)
        \item Data source: Hensel et al. \textit{Nature Physics} 2024 (simulated)
    \end{itemize}

\end{frame}

% ============================================
% SLIDE 5: Features
% ============================================
\begin{frame}{Input Features (15)}

    \begin{columns}[T]
        \begin{column}{0.48\textwidth}
            \textbf{Photon ratios (6)}
            \begin{itemize}
                \item $n_i / N_{total}$ for each position
                \item Normalized $\rightarrow$ brightness-invariant
            \end{itemize}

            \vspace{0.4cm}

            \textbf{Beam positions (6)}
            \begin{itemize}
                \item Normalized coordinates
                \item Where was the donut centered?
            \end{itemize}
        \end{column}

        \begin{column}{0.48\textwidth}
            \textbf{Modulation depth (2)}
            \begin{equation*}
                \text{mod}_x = n_{X^-} + n_{X^+} - 2 \cdot n_{X^0}
            \end{equation*}
            \begin{itemize}
                \item Measures signal contrast
                \item High $\rightarrow$ molecule off-center
            \end{itemize}

            \vspace{0.4cm}

            \textbf{Log total photons (1)}
            \begin{itemize}
                \item $\log(N_{total})$
                \item Proxy for signal-to-noise ratio
            \end{itemize}
        \end{column}
    \end{columns}

\end{frame}

% ============================================
% SLIDE 4: Results - Performance
% ============================================
\begin{frame}{Results}

    \begin{columns}[T]
        \begin{column}{0.48\textwidth}
            \textbf{Speed}
            \vspace{0.3cm}

            \begin{tabular}{lcc}
                \toprule
                Metric & MLE & XGBoost \\
                \midrule
                Inference Time & 1--10 ms* & 0.2 ms \\
                \textbf{Speedup} & 1$\times$ & \textbf{5--50$\times$} \\
                \bottomrule
            \end{tabular}

            \vspace{0.3cm}
            {\footnotesize *Literature estimate (Balzarotti et al. 2017)}
        \end{column}

        \begin{column}{0.48\textwidth}
            \textbf{Accuracy}
            \vspace{0.3cm}

            \begin{tabular}{lcc}
                \toprule
                Distance & RMSE & Bias \\
                \midrule
                15 nm & 1.7 nm & +0.8 nm \\
                20 nm & 3.1 nm & +1.3 nm \\
                30 nm & 4.2 nm & $-$2.6 nm \\
                \midrule
                \textbf{Overall} & \textbf{3.2 nm} & -- \\
                \bottomrule
            \end{tabular}
        \end{column}
    \end{columns}

    \vspace{0.5cm}

    \textbf{Uncertainty Quantification (MAPIE):}
    \begin{itemize}
        \item 90\% confidence intervals for each prediction
        \item Quantifies prediction reliability in real-time
    \end{itemize}

\end{frame}

% ============================================
% SLIDE 5: Throughput Comparison
% ============================================
\begin{frame}{Throughput Comparison}

    \centering

    \textbf{In the time MLE processes 1 distance...}

    \vspace{0.8cm}

    \begin{columns}[c]
        \begin{column}{0.45\textwidth}
            \centering
            {\Large\textbf{MLE}}

            \vspace{0.5cm}

            {\Huge 1}

            \vspace{0.3cm}

            {\footnotesize estimation}

            \vspace{0.3cm}

            {\small (10 ms)}
        \end{column}

        \begin{column}{0.1\textwidth}
            \centering
            {\Huge$\rightarrow$}
        \end{column}

        \begin{column}{0.45\textwidth}
            \centering
            {\Large\textbf{XGBoost}}

            \vspace{0.5cm}

            {\Huge 50}

            \vspace{0.3cm}

            {\footnotesize estimations}

            \vspace{0.3cm}

            {\small (0.2 ms each)}
        \end{column}
    \end{columns}

    \vspace{1cm}

    {\footnotesize Assuming 10 ms MLE inference time (upper bound estimate)}

\end{frame}

% ============================================
% SLIDE 6: Why It Works
% ============================================
\begin{frame}{Why It Works}

    \textbf{Precision follows Fisher Information:}
    \begin{itemize}
        \item 15 nm: Steeper PSF slope $\rightarrow$ higher Fisher Information $\rightarrow$ lower RMSE (1.7 nm)
        \item 30 nm: Flatter PSF slope $\rightarrow$ lower Fisher Information $\rightarrow$ higher RMSE (4.2 nm)
    \end{itemize}

    \vspace{0.5cm}

    \textbf{Consistent with Cramér-Rao Bound:}
    \begin{equation*}
        \sigma \propto \frac{L}{\sqrt{N}}
    \end{equation*}

    \vspace{0.5cm}

    \textbf{Key Insight:} Model learned MINFLUX physics without explicit constraints

\end{frame}

% ============================================
% SLIDE 7: Limitations
% ============================================
\begin{frame}{Limitations}
    
    \textbf{Training Data:}
    \begin{itemize}
        \item Only simulated data (Hensel et al. 2024)
        \item Limited range: 15, 20, 30 nm
        \item Single photon budget ($\sim$200 photons)
    \end{itemize}
    
    \vspace{0.5cm}
    
    \textbf{Validation:}
    \begin{itemize}
        \item No experimental validation yet
        \item Systematic bias at boundaries
        \item MLE timing based on literature estimates
    \end{itemize}
    
    \vspace{0.5cm}
    
    \textbf{Next Step:} Validation with experimental data
    
\end{frame}

% ============================================
% SLIDE 8: Future Work
% ============================================
\begin{frame}{Future Work}
    
    \begin{enumerate}
        \item \textbf{Experimental Validation}
        \begin{itemize}
            \item Test on real MINFLUX measurements
            \item Compare against MLE results
        \end{itemize}
        
        \vspace{0.3cm}
        
        \item \textbf{Pipeline Integration}
        \begin{itemize}
            \item Real-time acquisition system
            \item Live feedback during experiments
        \end{itemize}
        
        \vspace{0.3cm}
        
        \item \textbf{Extensions}
        \begin{itemize}
            \item 3D localization (currently 2D only)
            \item Multiple photon budgets
            \item Transfer learning to experimental data
        \end{itemize}
    \end{enumerate}
    
\end{frame}

% ============================================
% SLIDE 9: Summary
% ============================================
\begin{frame}{Summary}
    
    \begin{itemize}
        \item[\checkmark] \textbf{5--50$\times$ speedup} (1--10 ms $\rightarrow$ 0.2 ms)
        \item[\checkmark] \textbf{3.2 nm RMSE} (comparable accuracy)
        \item[\checkmark] \textbf{Learned physical principles} (Fisher Information)
        \item[\checkmark] \textbf{Uncertainty quantification} with MAPIE
    \end{itemize}
    
    \vspace{0.8cm}
    
    \centering
    {\Large Removes computational bottleneck for MINFLUX distance estimation}
    
    \vspace{1cm}
    
    {\Large\textbf{Thank you!}}
    
    \vspace{0.3cm}
    
    {\large Questions?}
    
\end{frame}

% ============================================
% BACKUP SLIDES
% ============================================
\appendix

\begin{frame}{Backup: Feature Engineering Details}
    
    \textbf{Photon ratios:}
    \begin{itemize}
        \item Normalized by total photons
        \item Invariant to brightness variations
    \end{itemize}
    
    \vspace{0.3cm}
    
    \textbf{Modulation depth:}
    \begin{equation*}
        \text{mod}_x = N(X^-) + N(X^+) - 2 \cdot N(X^0)
    \end{equation*}
    
    \vspace{0.3cm}
    
    \textbf{Normalized positions:}
    \begin{itemize}
        \item Uses training statistics (not per-sample!)
        \item Critical for generalization
    \end{itemize}
    
\end{frame}

\begin{frame}{Backup: Training Details}
    
    \textbf{Data:}
    \begin{itemize}
        \item Source: Hensel et al. \textit{Nature Physics} 2024
        \item Balanced sampling: 25k per distance (15/20/30 nm)
        \item Train/test split: 80/20
    \end{itemize}
    
    \vspace{0.5cm}
    
    \textbf{Hyperparameters:}
    \begin{itemize}
        \item n\_estimators: 500
        \item max\_depth: 6
        \item learning\_rate: 0.1
        \item subsample: 0.8
    \end{itemize}
    
    \vspace{0.3cm}
    
    \textbf{Training time:} $\sim$2 minutes on CPU
    
\end{frame}

\end{document}